% title: Prime Square Minus One Divisible By Twenty Four
% short-title: Prime Square Divisibility
% abstract: We prove the standard divisibility fact that if \(p\) is a prime greater than 3, then 24 divides \(p^2-1\). The proof combines the congruence of odd squares modulo 8 with the congruence of nonmultiples of 3 modulo 3.
% subjclass: 11A07, 11A41
% keywords: prime numbers, divisibility, congruences, Lean 4
% author: Mario Hernández M.

\subsection{Prime Square Divisibility}

\begin{problemstatement}
Prove that for any prime \(p > 3\), \(p^2-1\) is divisible by 24.
\end{problemstatement}

\begin{theorem}
If \(p\) is prime and \(p>3\), then
\[
  24 \mid p^2-1.
\]
\end{theorem}

\begin{proof}
Since \(p>3\), the prime \(p\) is neither \(2\) nor \(3\). In particular
\(p\) is odd. Every odd number has the form \(2k+1\), and then
\[
  (2k+1)^2-1=4k(k+1).
\]
The product \(k(k+1)\) of two consecutive integers is even, so this difference
is divisible by \(8\). This is formalized as
\lean{Biblioteca.Demonstrations.eight_dvd_square_sub_one_of_odd}.

Also \(3\nmid p\). A number not divisible by \(3\) is congruent to either
\(1\) or \(2\) modulo \(3\), and in both cases its square is congruent to
\(1\) modulo \(3\). Hence \(3\mid p^2-1\). This auxiliary congruence is
\lean{Biblioteca.Demonstrations.three_dvd_square_sub_one_of_not_dvd}.

Finally, \(8\) and \(3\) are coprime, so a number divisible by both is divisible
by \(8\cdot 3=24\). Applying this to \(p^2-1\) proves the theorem, recorded in
Lean as \lean{Biblioteca.Demonstrations.prime_square_sub_one_dvd_twenty_four}.
\end{proof}
